\documentclass[11pt]{article}
\usepackage[utf8]{inputenc}
\usepackage{amsmath,amssymb}
\usepackage[margin=1in]{geometry}
\usepackage{hyperref}
\emergencystretch=3em

\title{Weighted one-dimensional power substitution}
\author{Claude, with Daniel Murfet}
\date{2026-09-07}

\begin{document}
\maketitle
\sloppy

\begin{abstract}
Step~1 of the general-$d$ equal-ratio monomial milestone (Astra~\#14). For an arbitrary measurable
$g\colon\mathbb R\to[0,\infty]$,
\[
\int_{(0,1]}u^h\,g(u^{2k})\,du=\frac1{2k}\int_{(0,1]}t^{\lambda-1}g(t)\,dt,\qquad\lambda=\frac{h+1}{2k},
\]
in Lebesgue-integral form (\texttt{lintegral\_weighted\_power\_subst}), so no integrability is needed
and the identity can be iterated inside Tonelli arguments. It rests on the change of variables
$t=u^p$ on $(0,1]$ (\texttt{lintegral\_Ioc\_rpow\_subst}), which maps the interval onto itself
(\texttt{image\_rpow\_Ioc}), via Mathlib's one-dimensional Jacobian formula. Zero
\texttt{sorry}/\texttt{axiom}.
\end{abstract}

\section{The things formalised}
{\small
\begin{verbatim}
theorem image_rpow_Ioc (p) (hp : 0 < p) : (fun x : ℝ => x ^ p) '' Ioc 0 1 = Ioc 0 1
theorem lintegral_Ioc_rpow_subst (p) (hp : 0 < p) (g : ℝ → ENNReal) :
    ∫⁻ t in Ioc 0 1, g t = ∫⁻ u in Ioc 0 1, ENNReal.ofReal (p * u ^ (p - 1)) * g (u ^ p)
theorem lintegral_weighted_power_subst (h k : ℕ) (hk : 0 < k) (G : ℝ → ENNReal) :
    ∫⁻ u in Ioc 0 1, ENNReal.ofReal (u ^ h) * G (u ^ (2 * k))
      = ENNReal.ofReal (1 / (2 * k))
        * ∫⁻ t in Ioc 0 1, ENNReal.ofReal (t ^ ((h + 1) / (2 * k) - 1)) * G t
\end{verbatim}
}

\section{Lean notes}
\begin{itemize}
\item \texttt{MeasureTheory.lintegral\_image\_eq\_lintegral\_abs\_deriv\_mul} (file
\texttt{JacobianOneDim}) with $s=(0,1]$, $f=(\cdot)^p$, derivative from
\texttt{Real.hasDerivAt\_rpow\_const} and injectivity from \texttt{Real.rpow\_left\_injOn}
restricted to nonnegative reals.
\item Rewriting the substitution into the right-hand side only needs \texttt{conv\_rhs}; a bare
\texttt{rw} hits the left-hand integral first.
\item \texttt{ENNReal.ofReal\_mul} with metavariables rewrites the first product it finds; pass
the factors explicitly (\texttt{@ENNReal.ofReal\_mul (1/p) (p * u\^{}(p-1)) …}) before \texttt{ring},
which treats the \texttt{ofReal} atoms consistently only if both sides are split identically.
\end{itemize}

\end{document}
